\chapter{Conclusion}

This report documented CrisisSense, a frozen three-class crisis-language classification system that runs a character-level TF-IDF + Logistic Regression model and a BERT sequence-classification model independently on submitted text, without combining their outputs into an ensemble. Both models were evaluated on an identical, frozen 195-row test set drawn from a 1{,}294-row dataset split 905/194/195 into train/validation/test.

On this shared test set, BERT achieved higher overall accuracy (76.92\% vs.\ 71.28\%) and macro-F1 (74.66\% vs.\ 69.62\%) than TF-IDF + Logistic Regression, and a substantially higher F1 on the explicit-crisis class (84.66\% vs.\ 77.01\%). On the implicit-crisis class specifically, the two models were essentially tied, with TF-IDF marginally ahead (67.92\% vs.\ 67.33\% F1). BERT's mean selected-class confidence (94.60\%) was far higher than TF-IDF's (55.74\%), but this was accompanied by a higher Expected Calibration Error (0.177 vs.\ 0.155), not a lower one, indicating that BERT's higher confidence does not correspond to proportionally higher reliability on this test set. The two models agreed on 75.9\% of test predictions and disagreed on 24.1\%; agreement was shown to correlate with, but not guarantee, correctness for either model.

This project demonstrates that overall accuracy and macro-F1 do not tell the whole story of a model comparison: a model that leads clearly in aggregate metrics can still be matched, on a specific and practically important subclass, by a much simpler and cheaper baseline. It also demonstrates that a large gap in reported confidence between two models is not evidence that the more confident model is more trustworthy, once that confidence is checked against calibration error rather than taken at face value. CrisisSense classifies crisis-related language in text; it does not diagnose individuals, does not assess suicide risk, and its results in this report are scoped strictly to the frozen 195-sample evaluation set on which they were measured.
